\section{引言}
\label{sec:introduction}

功能材料的综合表征和标准化生物检测的执行——涵盖表面润湿性、耐腐蚀性、抗菌活性和生物反应性——对于推进生物医学工程、材料科学和生物技术至关重要。传统上，这些多阶段实验工作流程需要劳动密集型的手动实验室程序。例如，制备单个防腐薄膜样品涉及按精确试剂比例进行顺序溶液制备、在金属基底上进行受控表面涂覆，以及通过接触角测量或电化学测试进行后续表征。类似地，建立抗菌剂的最低抑菌浓度（MIC）分布需要在96孔微孔板中进行细致的梯度稀释，然后进行长时间培养孵育。这些手动工作不仅繁琐且通量有限，而且极易受到人为错误、交叉污染和液体溢出的影响，严重损害实验可重复性和操作安全性——特别是在处理腐蚀性或生物危害试剂时。

尽管机器人自动化已开始改变实验室操作，但其在复杂多工作流实验管道中的应用仍受到若干基本限制。首先，现有自动化工作站通常依赖刚性物理示教和预编程航点，虽然对简单重复运输有效，但缺乏异构工作流多样化操作需求所需的动态适应性——例如溶液制备期间的精密微量移液、薄膜制造期间的受控浸涂速度以及梯度稀释期间的微体积液体处理。其次，大多数部署的系统采用单个机械臂，当复杂工作流需要在多个物理分离的工作站进行操作时，限制了通量和空间覆盖范围。第三，跨长流程协调异构技能需要明确的完成判据与失败恢复机制。最后，在涉及腐蚀性薄膜前体或生物危害微生物培养物的物理实验室中直接调试硬编码的机器人轨迹，风险极高且成本昂贵。

为解决这些挑战，我们提出了一种面向生化实验室工作流程的多平台混合机器人框架，覆盖按受控试剂混合比例进行精密溶液制备、通过浸涂在金属基底上制备防腐薄膜，以及在96孔微孔板中进行微生物培养梯度稀释等任务。我们的系统架构由三个物理分离但协同工作的机器人平台组成（图~\ref{fig:overall_framework}）：

\textbf{Franka Panda臂（工作站A——溶液制备）：}一个配备关节力矩传感和阻抗控制的7自由度机械臂，负责抓取试剂瓶、精密容量移液、受控倾倒溶液混合，以及将制备好的溶液放置到移动机器人的运输托盘上。

\textbf{差速驱动移动机器人（站间运输）：}一个配备稳定优化运输托盘（带有防滑硅胶垫和ArUco基准标记）的自主移动平台，负责在制备站和分析站之间无溢出地运输开敞液体容器。

\textbf{UR3臂（工作站B——表面涂覆和生物检测）：}一个配备Robotiq 2F-85平行夹爪的紧凑型6自由度机械臂，负责从移动托盘取回容器、在金属基底上执行受控浸涂操作进行防腐薄膜制造，以及在96孔微孔板中执行自动梯度稀释进行微生物培养制备。

我们在NVIDIA Isaac Sim中进行并行策略训练，并在物理部署前测试操作技能。关键的是，我们提出了一种宏微解耦控制策略来处理异构任务需求。对于宏观站间运输，我们将移动机器人的导航表述为受约束的经典优化问题，采用最小加加速度轨迹优化和基于容器几何形状和填充比例的动态速度界，从根本上消除液体晃动和溢出风险。对于每个工作站的精密关键操作，我们实施平台特定的深度强化学习（DRL）和模仿学习（IL）管道：

对于溶液制备（Panda臂），我们训练姿态感知PPO智能体进行直立抓取，并实现基于IL的移液策略，通过正则化少样本真实世界微调实现微升级别的体积精度。对于薄膜制造（UR3臂），我们开发DRL策略来学习控制浸涂速度曲线和基底浸入角度，通过惩罚薄膜厚度变化的奖励函数优化涂层均匀性。对于梯度稀释（UR3臂），我们训练IL策略进行顺序多孔移液，具有自动更换吸头行为，确保一致的稀释比例同时防止孔间交叉污染。

工作流程中包含工作站间的容器转移。在任务级编排方面，我们采用具备失败恢复机制的分层状态机及具有标准化输入输出接口的可复用技能库。预定义条件根据实际任务观测决定技能切换与恢复动作。

总之，本工作包括以下系统模块及状态同步接口：
\begin{itemize}
    \item \textbf{仿真与实体状态接口：}我们定义机械臂关节、夹爪状态与物体位姿的时间戳同步方式，用于比较仿真与实体轨迹，其实测汇总结果见第\ref{subsec:state_sync_eval}节。
    \item \textbf{并行仿真训练：}我们在Isaac Sim中构建了一个统一的高保真仿真环境，涵盖机器人平台与实验室工作流程，支持无物理风险的并行策略训练。
    \item \textbf{异构工作流的宏微解耦混合控制：}我们提出了一种控制架构，将宏观移动机器人运输委托给约束感知的经典规划，同时利用平台特定的DRL和IL策略，针对化学合成、表面涂覆和生物制备的不同操作需求进行定制。
    \item \textbf{层次化技能编排：}我们构建了由具备失败恢复机制的分层状态机编排的可复用技能库，明确规定完成、超时及恢复条件。
\end{itemize}
